\section{Related Work}
\label{sec:related}

\paragraph{Human behaviour in ceremony and protocol verification.} Basin et al.~\cite{basin2016modeling} model untrained humans who may perform any action, and restrict this behaviour by rules. Bella et al.~\cite{bella2022modelling} model human threats as deviations of human agents. Fila and Radomirovi\'c~\cite{fila2024nothing} model out-of-band channels and information loss inside Tamarin, so forgetting enlarges the state the verifier explores. \XMen{}~\cite{sempreboni2020x,sempreboni2023mutation,quamara2025user} instead mutates the specification before verification, and \forget~\SCP{} keeps knowledge monotone. We follow this design, so the verifier analyses ordinary Tamarin theories. Models of human error in interactive systems~\cite{curzon2007approach} and robustness analyses against environmental deviations~\cite{zhang2023fortis} address related questions for other systems; none distinguishes lost values from lost skills or represents recall.

\paragraph{Defeasible reasoning.} We use DL($\partial$) with team defeat~\cite{nute1994defeasible,antoniou2001representation}; the individual-defeat and ambiguity-propagating variants are from~\cite{billington1993defeasible,antoniou2000flexible,governatori2004argumentation}, and propositional DL has a linear-time algorithm~\cite{maher2001propositional}, which our reasoner does not implement. Mature reasoners~\cite{lam2009spindle} and ASP encodings~\cite{governatori2026ddl} exist; we used the latter as an oracle for both DL($\partial$) and DL($\delta$). We compile a specific, analysable fragment, and our results concern how it depends on the variant of the logic. Forgetting in knowledge representation~\cite{lin1994forget,eiter2019brief} eliminates symbols from theories; it is not a model of human memory.

\paragraph{Mutation, validation and cognitive grounding.} \XMen{} mutants describe human deviations rather than the seeded faults of mutation testing~\cite{demillo1978hints,jia2011analysis}. The non-reuse invariant of RQ5 is per-model validation in the spirit of~\cite{pnueli1998translation}, weaker than a refinement proof (\cref{sec:discussion}). The distinctions we formalise come from the psychology of memory: knowing-how versus knowing-that~\cite{ryle1949concept,cohen1980preserved,anderson1982acquisition}, accessibility versus availability and cue-dependent recall~\cite{tulving1974cue,tulving1973encoding}, and reconstructive remembering~\cite{bartlett1995remembering,schacter1999seven}. The adversarial extension draws on misinformation and retrieval interference~\cite{loftus2005planting,anderson1994remembering};  these sources motivate modelling choices, not predictions of behaviour.